% ======================================================
% SLIDE: KẾ HOẠCH THỰC HIỆN – GIAI ĐOẠN 1
% ======================================================
\begin{frame}
    \frametitle{Kế hoạch thực hiện – Giai đoạn 1}

    \begin{BulletTight}
        \item Tìm hiểu về các nghiên cứu LLM và hệ đa tác tử, phân tích các hướng tiếp cận và công cụ kiểm thử hiện có.
        \item Xác định vấn đề, khoảng trống và mục tiêu nghiên cứu.
        \item Xác định các vai trò tác tử và cơ chế giao tiếp, phối hợp giữa các tác tử.
        \item Thiết kế kiến trúc tổng thể; chuẩn bị dữ liệu và baseline đánh giá.
        \item[\raise0.2ex\hbox{\color{KHTN_BLUE!85}\scriptsize\faArrowRight}]
        Kết quả: tổng quan nghiên cứu, đề xuất hệ thống, kế hoạch chi tiết;
        \textbf{01 bài báo hội nghị}.
    \end{BulletTight}
\end{frame}

% ======================================================
% SLIDE: KẾ HOẠCH THỰC HIỆN – GIAI ĐOẠN 2
% ======================================================
\begin{frame}
    \frametitle{Kế hoạch thực hiện – Giai đoạn 2}
    \begin{BulletTight}
        \item Phát triển các tác tử: Planner, API/UI, Evaluator, Coordinator, HITL\ldots
        \item Triển khai cơ chế giao tiếp, phối hợp giữa các tác tử và con người
        \item Đi sâu nghiên cứu một tác tử đại diện (Vd: API/UI)
        \item Tối ưu hiệu năng, ghi nhận lỗi, hạn chế và cải tiến.
        \item[\raise0.2ex\hbox{\color{KHTN_BLUE!85}\scriptsize\faArrowRight}]
        Kết quả: nguyên mẫu hệ thống; \textbf{02 bài báo hội nghị}.
    \end{BulletTight}
\end{frame}

% ======================================================
% SLIDE: KẾ HOẠCH THỰC HIỆN – GIAI ĐOẠN 3
% ======================================================
\begin{frame}
    \frametitle{Kế hoạch thực hiện – Giai đoạn 3}
    \begin{BulletTight}
        \item Mở rộng thực nghiệm; đo hiệu năng \& độ tin cậy.
        \item Tối ưu hóa cơ chế phối hợp, giao tiếp.
        \item Tối ưu quy trình HITL và giảm hallucination.
        \item Tinh chỉnh \textit{tác tử trọng tâm}
        \item Chuẩn hoá hệ thống; hoàn thiện luận án.
        \item[\raise0.2ex\hbox{\color{KHTN_BLUE!85}\scriptsize\faArrowRight}]
        Kết quả: hệ thống hoàn chỉnh; \textbf{01 bài báo tạp chí}.
    \end{BulletTight}
\end{frame}